\documentclass[11pt]{article}

\usepackage[margin=0.82in]{geometry}
\usepackage[T1]{fontenc}
\usepackage{lmodern}
\usepackage{microtype}
\usepackage{amsmath,amssymb,mathtools}
\usepackage{booktabs,longtable,tabularx,array,multirow}
\usepackage{pdflscape}
\usepackage[table]{xcolor}
\usepackage{enumitem}
\usepackage{siunitx}
\usepackage[hidelinks]{hyperref}
\usepackage{fancyhdr}

\definecolor{facetblue}{RGB}{32,98,155}
\definecolor{facetgreen}{RGB}{37,132,82}
\definecolor{facetred}{RGB}{190,62,54}
\definecolor{lightblue}{RGB}{235,243,249}
\definecolor{lightgray}{RGB}{245,245,245}

\hypersetup{
  pdftitle={Hadronic-Coefficient Simplification in FACET},
  pdfauthor={FACET calculation record}
}

\pagestyle{fancy}
\fancyhf{}
\lhead{FACET calculation record}
\rhead{Hadronic-coefficient simplification}
\cfoot{\thepage}
\setlength{\headheight}{14pt}

\setlist[itemize]{leftmargin=1.35em,itemsep=2pt,topsep=4pt}
\setlist[enumerate]{leftmargin=1.55em,itemsep=2pt,topsep=4pt}
\newcommand{\code}[1]{\texttt{#1}}
\newcommand{\eps}{\epsilon}
\newcommand{\GLI}{\mathrm{GLI}}
\newcommand{\bytes}{\,\mathrm{bytes}}
\newcommand{\Sec}{\,\mathrm{s}}
\newcommand{\good}{\textcolor{facetgreen}{\textbf{useful}}}
\newcommand{\bad}{\textcolor{facetred}{\textbf{not retained}}}

\title{\vspace{-1.1cm}\textbf{Hadronic-Coefficient Simplification in FACET}\\[0.35em]
\large A measured NLO and NNLO comparison of algebraic orderings}
\author{Internal calculation record}
\date{9 August 2026}

\begin{document}
\maketitle

\begin{abstract}
This report reconstructs the hadronic-coefficient simplification calculations
performed during 7--8 August 2026 and records their consolidation into a
process-card-driven FACET routine.  It distinguishes algebra done on
amplitude-pair coefficients, complete unreduced integral coefficients, and
final master-integral coefficients.  Every quoted size is identified as either
Mathematica \code{ByteCount} or an actual serialized file size.  At NLO,
simplifying complete target coefficients before inserting the Kira reduction
rules is the best measured ordering; a bounded master-level cleanup is then
inexpensive.  At NNLO, constructing a monolithic master coefficient is not
viable.  The measured alternative retains additive contributions, performs
exact rational algebra in positive momentum-fraction roots, and certifies the
common Laurent monomial only after cancellations.  NLO and NNLO now use one
physical normalization declared by the process card; they differ only in
whether a coefficient is supplied already assembled or as additive
contributions.  The reviewed implementation separately certifies the source
kinematic chamber and its dimensionless image, verifies the mass dimension and
invertibility of the coordinate map, and freezes branch-sensitive analytic
objects while rational coefficient algebra is performed.  The report also
isolates the TT BMHV dimension-bookkeeping
error and separates coefficient algebra from the still-open analytic
evaluation of NNLO cut masters.
\end{abstract}

\tableofcontents
\clearpage

\section{Scientific object and scope}

The object being reduced is an exact coefficient multiplying a cut master
integral.  For an amplitude labelled by $p$ and its conjugate labelled by $q$,
write
\begin{align}
  \mathcal A_p\mathcal A_q^{*}
  &=\sum_{T} c_{T}^{(pq)}\,T,
  &
  T&=\sum_m R_{Tm}\,M_m .
\end{align}
Here $T$ is a cut integral before IBP reduction, $M_m$ is a master integral,
and $R_{Tm}$ is the exact Kira reduction coefficient.  The final coefficient
of $M_m$ is therefore
\begin{equation}
  C_m=\sum_{p,q,T}c_T^{(pq)}R_{Tm}.
  \label{eq:master-coeff}
\end{equation}
This identity defines three distinct places where algebra can be attempted:
\begin{enumerate}
  \item simplify each $c_T^{(pq)}$ before summing amplitude pairs;
  \item form $c_T=\sum_{p,q}c_T^{(pq)}$, simplify it, and then insert
        $R_{Tm}$;
  \item form the complete $C_m$ and simplify only after IBP substitution.
\end{enumerate}

Let the nontrivial momentum fractions be
$\boldsymbol{\xi}=(\xi_1,\ldots,\xi_r)$.  After exact hadronic-coordinate
substitution, every nonzero coefficient accepted by the simplifier must have
the form
\begin{equation}
  C_m=\mathcal D(\boldsymbol{\xi})
  \boldsymbol{\xi}^{\boldsymbol{\nu}} H_m(I,\eps),
  \qquad
  \boldsymbol{\xi}^{\boldsymbol{\nu}}
  \equiv\prod_{i=1}^{r}\xi_i^{\nu_i},
  \label{eq:generic-hard-factor}
\end{equation}
where $\mathcal D$ is the product of twist-two distributions, $I$ denotes
the chosen partonic invariants, and $H_m$ contains neither momentum fractions
nor forbidden hadronic vectors.  The card may state $\mathcal D$ and
$\boldsymbol{\nu}$ explicitly, or request their exact derivation.

For the ppHX UU example this gives
\begin{equation}
  \mathcal P_{\rm UU}=
  \frac{f_1(x_a)f_1(x_b)D_1(z_h)}{x_a x_b z_h^2},
  \qquad
  C_m=\mathcal P_{\rm UU}\,H_m(s,t,u,\eps),
  \label{eq:universal-factor}
\end{equation}
up to coupling, color, and conventional powers of $\pi$.  Equation
\eqref{eq:universal-factor} is an example supplied by the ppHX card, not an
assumption embedded in the reusable simplifier.

This report concerns coefficient algebra.  It does not claim that the 342
NNLO cut masters have been evaluated analytically, nor that endpoint and
distributional reconstruction has been completed.

\section{What every quoted size means}

Several earlier notes used the word ``bytes'' for different objects.  The
distinction is essential:
\begin{description}[leftmargin=3.1cm,style=nextline]
  \item[Mathematica \code{ByteCount}] An in-kernel estimate of the expression
  tree.  It is meaningful for comparisons in the same Mathematica version.
  \item[Serialized bytes] The exact \code{FileByteCount} of a saved \code{.wl}
  or binary artifact.  It measures I/O and disk use and depends on the chosen
  serialization grammar.
  \item[Additive terms] The number of summands after splitting only the
  top-level \code{Plus}.  A different outer factorization changes this count.
  \item[Objects] Independently scheduled pair, target, or master
  coefficients.  Counts at different algebraic stages are not comparable.
  \item[Target leaves] Individual Kira-rule contributions entering a selected
  master column.
  \item[Fractions] Stored exact numerator-denominator pairs after local
  hadronic cleanup.
\end{description}
No reduction ratio in this report mixes \code{ByteCount} with serialized
bytes.

\section{Physical substitution and branch control}

\subsection{The process-card contract}

Every card used for hard-coefficient simplification contains a
\code{CoefficientKinematics} block declaring
\begin{enumerate}
  \item the positive fractions $\boldsymbol{\xi}$, either explicitly or by
        exact inference from the card's momentum-fraction map;
  \item the expected distribution product $\mathcal D$ and Laurent valuation
        $\boldsymbol{\nu}$, or \code{Automatic} for exact derivation;
  \item the physical invariant chamber;
  \item a positive scale $Q^2$ and an invertible map to dimensionless
        invariant coordinates;
  \item variables forbidden in $H_m$, with automatic inclusion of all
        declared fractions and hadronic basis vectors;
  \item the allowed branch grammar.
\end{enumerate}
The routine rejects a missing block, a missing declaration, a noninvertible
coordinate map, a dimensionful coordinate, a scale whose positivity is not
implied by the card, and an explicit fraction list that differs from the
momentum-fraction map.  It also requires the source chamber and the
dimensionless chamber to be separately nonempty.  If
\(X_i=X_i(I_1,\ldots,I_n)\) denotes the forward coordinate map and
\(I_a=I_a(Q^2,X_1,\ldots,X_n)\) its inverse, the exact checks are
\begin{align}
 X_i\bigl(I(Q^2,X)\bigr)&=X_i,
 &
 I_a\bigl(Q^2,X(I)\bigr)&=I_a,\\
 \mathcal R_I&\Longrightarrow \mathcal R_X\bigl(X(I)\bigr),
 &
 \mathcal R_X&\Longrightarrow \mathcal R_I\bigl(I(Q^2,X)\bigr),
 \label{eq:chamber-map}
\end{align}
where \(\mathcal R_I\) and \(\mathcal R_X\) are the source and dimensionless
physical regions.  They are never conjoined before the coordinate equations
are imposed; this prevents an inconsistent card from making an equality check
vacuously true.

For ppHX, the NLO UU and TT cards contain exact coordinates for
\begin{equation}
  P_a,\ P_b,\ P_h,\ n_h,\ \bar n_h,\ S_T,\ S_{hT}
\end{equation}
in terms of $s,t,u,x_a,x_b,z_h$ and the spin azimuths.  The physical chamber
contains
\begin{equation}
  s>0,\qquad t<0,\qquad u<0,\qquad s+t+u>0,
  \qquad 0<x_a,x_b,z_h<1,
  \label{eq:physical-chamber}
\end{equation}
together with the required reality assumptions.  The package infers
$0<\xi<1$ for each nonempty momentum fraction declared in the card and
combines those inequalities with the card assumptions.

\subsection{Why known signs do not remove branch bookkeeping}

There is no physical ambiguity about the signs in
Eq.~\eqref{eq:physical-chamber}.  The algebraic issue is that Mathematica does
not globally replace, for example,
\begin{equation}
  \sqrt{x_a x_b}\longrightarrow \sqrt{x_a}\sqrt{x_b}
\end{equation}
unless positivity is built into the transformation.  The identity is false
on a general complex domain.  A global \code{PowerExpand} would therefore
erase information that the final analytic result must retain.

The implemented branch grammar, \code{PositiveMonomialRoots}, introduces one
positive root $r_i$ for each declared fraction,
\begin{equation}
  \xi_i=r_i^2,\qquad r_i>0,\qquad i=1,\ldots,r.
  \label{eq:root-lift}
\end{equation}
Only a certified Laurent monomial in the declared fractions is lifted.  Thus
\begin{equation}
  \prod_{i=1}^{r}\xi_i^{n_i/2}
  \longmapsto \prod_{i=1}^{r}r_i^{n_i}.
\end{equation}
A nonmonomial fraction base such as $\xi_1+\xi_2$ is rejected.  A logarithm,
Gamma function, hypergeometric function, absolute value, step function, or
distribution whose argument depends on a fraction is also rejected rather
than transformed.  Fraction-independent noninteger powers, logarithms, Gamma
functions, angular functions, distributions, and BMHV tensors are replaced by
inert symbols while the rational coefficient is simplified and are restored
afterward.  Exact equality is checked before and after this freezing; no
\code{PowerExpand} is used.

For an additive coefficient $C_m=\mathcal D\sum_j c_{mj}$, write each lifted
contribution as $c_{mj}=N_{mj}(\boldsymbol r)/D_{mj}(\boldsymbol r)$ and form
the exact sum
\begin{equation}
  \sum_j c_{mj}=\frac{N_m(\boldsymbol r)}{D_m(\boldsymbol r)}.
\end{equation}
The Laurent valuation is accepted only if there is an integer vector
$\boldsymbol n$ and a fraction-independent $H_m$ such that
\begin{equation}
  N_m(\boldsymbol r)=
  H_m\,\boldsymbol r^{\boldsymbol n}D_m(\boldsymbol r),
  \qquad
  \boldsymbol\nu=\frac{1}{2}\boldsymbol n.
  \label{eq:root-certificate}
\end{equation}
The identity in Eq.~\eqref{eq:root-certificate} is checked coefficient by
coefficient in the sparse root polynomial.  This is why square-root terms may
cancel between contributions without being discarded prematurely.

The scale and coordinates are likewise card data.  In the ppHX cards they are
\begin{equation}
  x=-\frac{t}{s},\qquad y=-\frac{u}{s},
  \qquad s>0,\quad x>0,\quad y>0,\quad x+y<1,
  \label{eq:dimensionless}
\end{equation}
and the routine derives the power $p_m$ from exact homogeneity:
\begin{equation}
  H_m(s,t,u,\eps)=s^{p_m}\widehat H_m(x,y,\eps).
\end{equation}
No common value of $p_m$ is assumed across masters, and $p_m$ may depend on
$\eps$.  The inverse coordinate map and the reconstruction of the dimensional
coefficient are checked exactly.

\section{NLO UU: where simplification should occur}

\subsection{Complete 5 by 5 calculation}

The complete 25-pair calculation gives six masters.  Table~\ref{tab:nlo5}
records the measured stages.

\begin{table}[ht]
\centering
\caption{NLO UU 5 by 5 calculation.  The final size is an actual serialized
file size; intermediate input sizes were not retained in this compact
summary.}
\label{tab:nlo5}
\begin{tabular}{@{}lr@{}}
\toprule
Quantity & Measured result \\
\midrule
Pair construction & $13.3465\Sec$ \\
Pair store & $1{,}086{,}418$ serialized bytes \\
Kira reduction & $50.1824\Sec$ \\
Unreduced targets & $77$ \\
Masters & $6$ \\
Hadronic substitution & $2.01\Sec$ \\
Complete-target cleanup & $6.65\Sec$ \\
Complete-master cleanup & $9.04\Sec$ \\
Complete coefficient stage & $48.5528\Sec$ \\
Final result & $170{,}483$ serialized bytes \\
\bottomrule
\end{tabular}
\end{table}

The extracted prefactor was
\begin{equation}
 \frac{C_F\,\pi^{3+\eps}\alpha_s^3
 D_1(z_h)f_1(x_a)f_1(x_b)}{x_a x_b z_h^2}.
\end{equation}
All six remaining coefficients were exactly free of momentum fractions and
hadronic basis vectors.

\subsection{The four orderings on the 10 by 10 sample}

The four measured orderings are summarized in Table~\ref{tab:nlo-order}.

\begin{table}[ht]
\centering
\caption{NLO UU orderings.  All sizes are Mathematica \code{ByteCount}.
The target-then-master time is the second stage only.}
\label{tab:nlo-order}
\scriptsize
\begin{tabular}{@{}lrrrrl@{}}
\toprule
Ordering & Objects & Input & Output & Time & Result after composition \\
\midrule
Pair first & 1332 & 89,111,000 & 2,694,944 & 43.403 s &
Master size 1,574,064 \\
Target first & 117 & 46,940,920 & 1,362,792 & 23.715 s &
Master size 1,548,952 \\
Master first & 1 tested & -- & -- & $>281$ s &
One coefficient unfinished \\
Target then master & 8 & 1,646,272 & 1,036,008 & 62.125 s &
Final size 1,041,912 \\
\bottomrule
\end{tabular}
\end{table}

Target first is the best first expensive operation because pair-to-pair
cancellations have already occurred while the Kira map has not yet inflated
the expression.  Pair first repeats algebra.  Master first presents the
simplifier with the largest sum.  Target then master gives the smallest final
object, but the sequential time is
\begin{equation}
  23.715\Sec+62.125\Sec=85.840\Sec,
\end{equation}
not 62.125 s as stated in the earlier narrative.

The compact benchmark summary did not finish its independent equality test
between all four routes.  The timing and size comparison is valid; exactness
claims in the final NLO calculation instead come from the later production
calculation and the certified fraction-normalization tests.

\subsection{Detailed NLO algebraic comparisons}

The largest target illustrates why a generic rational normal form is not
automatically a compact physics form:
\begin{center}
\begin{tabular}{@{}lrrr@{}}
\toprule
Method & Input \code{ByteCount} & Output \code{ByteCount} & Time \\
\midrule
Termwise \code{Simplify} & 3,517,376 & 49,016 & 4.838 s \\
Whole \code{Simplify} & 3,517,376 & 327,568 & 86.378 s \\
Whole \code{FullSimplify} & 3,517,376 & no result & 120.045 s \\
\code{FactorTerms[Cancel[Together]]} & 3,517,376 & 6,812,872 & 50.350 s \\
\bottomrule
\end{tabular}
\end{center}
\par\medskip
The isolated equality tests for the first two entries reached their time
bounds, so those two rows establish timing and representation size, not a
separate exact identity.  The complete later route was checked exactly.

The best NLO factor-first chain was:
\begin{align}
  1{,}362{,}792&\xrightarrow[3.393\Sec]{\text{extract common factor}}
  684{,}744,\\
  855{,}104&\xrightarrow[10.741\Sec]{\text{final master cleanup}}
  759{,}840,
\end{align}
where all sizes are \code{ByteCount}.  Including the preceding target cleanup
gives about 37.76 s.  Repeating the master cleanup for 9.426 s left the size
unchanged.

Direct target cleanup without factor extraction needed 450.188 s to change
1,542,536 to 727,376 bytes.  By contrast, certified division by the universal
fraction monomial followed by \code{Simplify} changed 1,035,984 to 283,144
bytes in 5.877 s.  A further \code{FullSimplify} gave only a modest reduction
to 256,456 bytes in 32.352 s.  \code{Factor}, color collection, signature
grouping, and denominator grouping all expanded the expression.  The complete
measurements appear in Appendix~\ref{app:all-methods}.

\section{NLO TT: correction, branch structure, and timing}

\subsection{The BMHV error}

TT contains evanescent tensor monomials.  Let a loop-independent factor from
the original trace be $F(D)$, and let $\mathcal T_r^{(D)}$ denote the rank
$2r$ evanescent tensor integral.  The correct bookkeeping is schematically
\begin{equation}
  F(D)\,\mathcal T_r^{(D)}
  \longrightarrow
  F(D)\sum_j a_j(D+2r)I_j^{(D+2r)},
  \label{eq:tt-correct}
\end{equation}
followed by dimensional recurrence for the scalar integrals.  The incorrect
code instead changed the inherited factor to $F(D+2r)$.

The diagnostic rank-two angular identity is
\begin{equation}
  \left\langle (k\cdot Y)^2\right\rangle_N
  =-\frac{k_\perp^2}{N-3},\qquad Y^2=-1.
  \label{eq:tt-toy}
\end{equation}
At tensor stage zero, $N=D$ and the denominator is $D-3$.  After one
dimension shift, $N=D+2$ and the denominator is $D-1$.  A pre-existing
loop-free factor such as $D-4$ remains $D-4$.  Exact tests checked rank one,
rank two, cut projection, both denominators, preservation of the original
factor, and rejection of the incorrect denominator.

The practical consequence was large:
\begin{center}
\begin{tabular}{@{}lrrr@{}}
\toprule
Route & \code{ByteCount} & GLI occurrences & Distinct GLIs \\
\midrule
Incorrect degree-zero tensor route & 184,796,536 & 39,036 & 31 \\
Correct sparse cut route & 27,754,864 & 9 & 9 \\
\bottomrule
\end{tabular}
\end{center}
\par\medskip
The TT code therefore did require a physics correction before simplification
measurements were meaningful.

\subsection{TT coefficient algebra}

For the 5 by 5 sample, target cleanup changed
139,298,664 to 10,562,384 \code{ByteCount} bytes in 52.02 s on four kernels.
Final-master cleanup took 30.36 s.  The final six-master file occupies 750,504
serialized bytes.  The angular dependence closes on
\begin{equation}
  \left\{\cos\phi_a\cos\phi_h,\ \sin\phi_a\sin\phi_h\right\},
\end{equation}
and both mixed sine-cosine structures vanish exactly.

For four representative TT targets, \code{Together} was fastest but returned
the larger rational form.  \code{Simplify}, \code{Together} followed by
\code{Simplify}, and the historical collect routine all reached the same
smaller size.  Reversing the order, \code{Simplify} followed by
\code{Together}, returned the larger form.  For example, the 6,438,104-byte
target changed to 22,848 bytes in 1.528 s with \code{Simplify}, to the same
size in 0.836 s with \code{Together} then \code{Simplify}, but to 74,048 bytes
in 0.480 s with the reverse order.

An explicit overall $i$ in an intermediate TT coefficient is not, by itself,
evidence of a wrong square-root branch.  It may originate from the amplitude,
cut, or projector convention.  Reality must be checked only after the full
normalization and angular projection.  The positive-root rules address the
separate issue of algebraically valid radical identities.

\section{NNLO: why monolithic master simplification fails}

The NNLO double-real UU reduction contains 44,877 Kira targets and 342
masters.  The earlier post-IBP master files entering the old final cleanup
occupy 3,960,695,102 serialized bytes; the resulting files occupy
3,597,827,125 bytes.  This is only a 1.101-fold reduction.

\subsection{Target-level tests}

The 64 largest target coefficients changed from 83,386,848 to 28,279,424
\code{ByteCount} bytes in 75.643 s on four kernels.  On the largest isolated
target, whole \code{Simplify} gave
\begin{equation}
  2{,}434{,}656\to811{,}576\quad(15.30\Sec),
\end{equation}
while termwise \code{Simplify} gave
\begin{equation}
  2{,}434{,}656\to701{,}328\quad(6.92\Sec).
\end{equation}
Both isolated differences were exactly zero.  Pure rational normalization was
faster but expanded the target: \code{Cancel[Together]} produced 3,748,912
bytes, and adding \code{FactorTerms} produced 5,316,312 bytes.

A stratified set of 125 complete targets changed from 52,630,888 to 559,984
\code{ByteCount} bytes in 30.45 s on four kernels.  The fraction dependence
consisted only of explicit half-integer Laurent powers; no
fraction-dependent logarithm, Gamma function, or hypergeometric function was
found in that sample.

\subsection{Final-master tests}

Eight whole-master \code{Simplify} jobs each reached 1,800 s and 4 GiB without
producing retained output.  Aggregate memory was approximately 34 GiB.
Termwise simplification of master 8 was feasible: 377 terms changed a
4,655,210-byte serialized file to 3,082,976 bytes in 175.070 s on eight
kernels.

The rational/signature decomposition gave only moderate serialized
reductions:
\begin{center}
\begin{tabular}{@{}lrrr@{}}
\toprule
Master & Input & Output & Time \\
\midrule
1 & 1,607,146,927 & 920,945,177 & about 19.5 min \\
2 & 903,734,615 & 464,281,765 & about 15 min \\
4 & 83,696,010 & 47,987,161 & 68.55 s \\
\bottomrule
\end{tabular}
\end{center}
\par\medskip
Master 4 had 2,591 terms and 2,591 different signatures, so there was no
cross-term signature merge.  The gain came only from rational normalization
inside individual terms.

Symbolica did not improve the tested coefficient: direct Mathematica took
5.91 s and produced 0.70 MB, while Symbolica followed by Mathematica took
12.16 s and produced 0.81 MB before including preprocessing time.  No exact
end-to-end FORM or SymPy route was established for the full grammar.

\section{NNLO: exact entrywise and denominator algebra}

\subsection{Entrywise positive-root cleanup}

Three post-IBP master columns were selected by incoming target-leaf count:
one small, one median, and the largest column.  Together they contain 38,460
target leaves in 38,400 exact denominator entries.  The source occupies
859,107,253 serialized bytes.  Applying the root lift
\eqref{eq:root-lift}, proving the universal Laurent factor entry by entry, and
reducing in the positive-$s$ root ring produced a 131,787,741-byte store in
1,304.639 s on eight kernels, a 6.52-fold reduction.

The independent audit checked source hashes, numerator and denominator hashes,
entry and leaf counts, exact arithmetic, nonzero denominators, absence of
momentum fractions and temporary roots, and integer powers of $s$.

\subsection{Equal-denominator assembly}

For a class of exactly equal denominators,
\begin{equation}
  \sum_i\frac{n_i}{d}=\frac{\sum_i n_i}{d}.
  \label{eq:equal-denom}
\end{equation}
A hash only locates candidate classes; structural equality decides whether
Eq.~\eqref{eq:equal-denom} is used.

Always retaining the cancelled form was a mistake for storage.  For the
largest master, the fraction count fell from 38,306 to 6,290 in 6,375.22 s,
but the fraction-list \code{ByteCount} grew from 422,282,520 to
1,727,567,608 and the serialized output occupied 488,886,853 bytes.

The corrected size-monotone rule constructs
\begin{align}
  \text{raw}&=\left\{\sum_i n_i,d\right\},\\
  \text{cancelled}&=\operatorname{NumDen}
  \!\left[\operatorname{Cancel}\!\left(\frac{\sum_i n_i}{d}\right)\right]
\end{align}
and retains the smaller exact pair according to \code{ByteCount}.  The largest
master then changed from 38,306 to 6,557 fractions in 1,260.216 s.  Its
fraction-list size changed from 422,282,520 to 351,559,840 bytes and its
serialized output occupied 99,892,985 bytes.

\subsection{Dimensionless normalization}

The exact homogeneity transformation in Eq.~\eqref{eq:dimensionless} found
scale powers $p_m=4,2,-2$ for the three selected masters.  For the largest
master it changed 351,559,840 to 296,802,344 \code{ByteCount} bytes and
99,892,985 to 92,385,524 serialized bytes in 430.914 s on eight kernels.

A second size-monotone denominator merge changed 6,557 to 5,973 fractions and
296,802,344 to 296,275,768 \code{ByteCount} bytes in 904.575 s.  The final
serialized file occupies 92,241,826 bytes.  A complete independent replay took
943.22 s.

Factoring every dimensionless denominator over $\mathbb Q(\eps)$ took 27.94 s
and found only 30 distinct $x,y$-dependent factors.  However, a separate
factor dictionary occupied 91,385,453 serialized bytes, only 0.93 percent less
than the direct fraction list.  It is useful for identifying the singular
locus, but not as a new storage representation.

\section{Is the simplification post-IBP?}

The answer depends on the row being discussed:
\begin{itemize}
  \item Pair-first and target-first simplification act before the Kira rules
  are inserted into Eq.~\eqref{eq:master-coeff}, although the rules have
  already been computed.
  \item Master-first, target-then-master's second stage, every final-master
  comparison from 7 August, and the complete NNLO selected-column root-ring
  study act after the IBP rules have been inserted.
  \item The NNLO selected-column entries are not unreduced integrals.  They are
  the individual post-IBP contributions $c_T R_{Tm}$ to a fixed master
  coefficient $C_m$.
\end{itemize}
Thus the measurements used to estimate the final 342 coefficients are
post-IBP master-coefficient algebra, as expected.

\section{Card-driven implementation and exact tests}

The reusable implementation is now contained in
\code{FeynFacet/Private/Simplification.wl}.  It contains no literal ppHX
fraction names, distribution product, Laurent valuation, invariant ratios, or
selected master labels.  These enter through the card contract stated above.
Historical benchmark drivers still contain ppHX-specific choices because they
record the measurements in this report; they are not the maintained
simplification interface.

The public routine \code{SimplifyHardCoefficients} applies the same physical
normalization to two exact input forms:
\begin{align}
  \text{assembled:}&\qquad \{C_1,\ldots,C_N\},\nonumber\\
  \text{contribution-wise:}&\qquad
  \bigl\{\{c_{11},\ldots,c_{1n_1}\},\ldots,
  \{c_{N1},\ldots,c_{Nn_N}\}\bigr\},
  \quad C_m=\sum_j c_{mj}.
  \label{eq:two-input-forms}
\end{align}
The second form is needed when constructing $C_m$ first would exceed memory or
when the cancellation in Eq.~\eqref{eq:root-certificate} must be exposed
without a global simplification.  It is not an NNLO-specific definition.
Either form may be used at any perturbative order.

The acceptance criteria in the exact tests were: the card schema is complete;
the source and dimensionless chambers satisfy Eq.~\eqref{eq:chamber-map}; every
coordinate is dimensionless; the distribution product and Laurent valuation
are derived or agree with their explicit declarations; no forbidden variable
remains; every scale-separated coefficient reconstructs the dimensional
expression exactly; and every result contains exact data only.  A renamed
synthetic process used
$\{\rho_A,\rho_B,\rho_H\}$, scale $Q^2$, and coordinates $X,Y$.  It also
included two contributions proportional to $\sqrt{\rho_A\rho_B}$ that cancel
only after addition.  The routine derived
\begin{equation}
  \boldsymbol\nu=(-1,-1,-2)
\end{equation}
and reconstructed both hard coefficients exactly.  The negative tests reject
a missing or incomplete block, contradictory coordinate chamber, dimensionful
coordinate, incomplete coordinate map, wrong twist-two channel, reciprocal
distribution factor, wrong Laurent valuation, residual parton momentum,
nonmonomial fraction root, and fraction-dependent analytic function.

Stored physics expressions gave the following measured results:
\begin{center}
\begin{tabular}{@{}lrrr@{}}
\toprule
Input & Exact coefficients & Time & Result \code{ByteCount} \\
\midrule
NLO UU & 6 & 48.51 s & 859.88 KiB \\
NLO TT, smallest retained coefficient & 1 & 2.58 s & 67.16 KiB \\
NNLO UU, one stored target & 1 (2 contributions) & 2.34 s & 27.91 KiB \\
\bottomrule
\end{tabular}
\end{center}
\par\medskip
For all three, the extracted distribution factor agreed with the card value.
The NLO cards declare and verify
\begin{equation}
  \mathcal D_{\rm extracted}=\mathcal D_{\rm card}
  =f_1(x_a)f_1(x_b)D_1(z_h),
  \qquad \boldsymbol\nu=(-1,-1,-2),
\end{equation}
and the extracted hard coefficient was exactly independent of
$x_a,x_b,z_h$.  The NNLO card currently derives the valuation from the exact
contribution sum; it should be made explicit only after the complete
342-master data establish a common value.  The NNLO timing is a test of one
two-contribution stored target, not an estimate for all 342 masters.

An independent end-to-end NLO calculation generated 16 cut integrals, reduced
them to one master, and reconstructed the final expression exactly.  Kira took
17.50 s and the coefficient stage took 10.34 s; the compact result occupied
122.48 KiB.  This test exercises the same public simplification routine used
for the stored-data calculations.

\section{Eight-kernel time estimate}

The selected post-IBP source occupies 859,107,253 serialized bytes.  The old
all-master post-IBP source occupies 3,960,695,102 bytes, giving the measured
size ratio
\begin{equation}
  \rho=\frac{3{,}960{,}695{,}102}{859{,}107{,}253}=4.61.
\end{equation}
Scaling the selected-column measurements by $\rho$ gives
\begin{align}
  t_{\rm roots}&\simeq 4.61\times1{,}304.64\Sec=1.67\ \mathrm{h},\\
  t_{\rm assemble+scale}&\simeq
  4.61\times(1{,}260.22+430.91+904.57)\Sec=3.32\ \mathrm{h}.
\end{align}
Uneven column sizes, serialization, scheduling, and exact audits motivate a
planning interval of
\begin{equation}
  \boxed{6\text{--}10\ \mathrm{hours}}
\end{equation}
for the complete post-IBP coefficient stage on eight kernels, starting from
the existing Kira reduction.  This extrapolation assumes that the selected
columns are representative per serialized source byte; it has not yet been
measured on all 342 columns.

A contemporaneous note quoted 1,313.82 s for one Kira calculation, but the
matching timing log was not found during this audit.  That number is not used
in the estimate above.

There is no defensible time estimate for the \emph{complete exact NNLO hard
function}.  Analytic evaluation of the 342 cut masters, boundary data,
endpoint expansion, plus distributions, and combination with the remaining
NNLO sectors are outside this one-day timing sample.  A numerical master check
cannot replace those analytic steps.

\section{Unified simplification route}

The physics sequence is common to NLO and NNLO:
\begin{enumerate}
  \item read the complete kinematic, fraction, distribution, scale, forbidden
        variable, and branch declarations from the card;
  \item substitute the exact hadronic coordinates;
  \item derive or verify the common twist-two distribution product;
  \item lift only certified half-integer Laurent monomials in positive
        fractions and prove Eq.~\eqref{eq:root-certificate};
  \item derive or verify the common fraction valuation and reject any remaining
        forbidden variable;
  \item derive each scale power from exact homogeneity, transform to the card's
        dimensionless invariants, and verify exact reconstruction.
\end{enumerate}
Both input forms are converted to contribution groups before step 2 and then
enter the same algebraic kernel.  Consequently perturbative order does not
select a different mathematical normal form.  The only scheduling difference
is whether contributions are summed before entering the routine or retained as
a bounded list until their distribution and Laurent factors have been
certified.
When complete $C_m$ expressions remain small, the assembled input in
Eq.~\eqref{eq:two-input-forms} is shorter and allows Mathematica to simplify
each coefficient directly.  When assembly is the memory bottleneck, the
contribution-wise input performs bounded exact denominator algebra before the
same physical checks.  This scheduling decision should be made from measured
expression size, not from the labels NLO and NNLO.

\clearpage
\appendix
\begin{landscape}
\section{Complete measured method inventory}
\label{app:all-methods}

The following tables include every method for which a numeric result was
retained.  A dash means that the corresponding metric was not measured or no
result was produced.

\scriptsize
\begin{longtable}{@{}llrrrl@{}}
\caption{NLO UU simplification measurements.}\label{tab:all-nlouu}\\
\toprule
Method & Object and metric & Input & Output & Time & Interpretation \\
\midrule
\endfirsthead
\toprule
Method & Object and metric & Input & Output & Time & Interpretation \\
\midrule
\endhead
Pair first & 1332 objects, \code{ByteCount} & 89,111,000 & 2,694,944 & 43.403 s & Repeats algebra; final master 1,574,064 bytes.\\
Target first & 117 objects, \code{ByteCount} & 46,940,920 & 1,362,792 & 23.715 s & Fastest first nontrivial route; final master 1,548,952 bytes.\\
Master first & one coefficient & -- & -- & $>281$ s & Unfinished.\\
Target then master & 8 objects, \code{ByteCount} & 1,646,272 & 1,036,008 & 62.125 s & Second stage only; sequential total 85.840 s.\\
Largest target, termwise \code{Simplify} & \code{ByteCount} & 3,517,376 & 49,016 & 4.838 s & Equality check later reached its time bound.\\
Largest target, whole \code{Simplify} & \code{ByteCount} & 3,517,376 & 327,568 & 86.378 s & Slower and larger than termwise.\\
Largest target, whole \code{FullSimplify} & \code{ByteCount} & 3,517,376 & -- & 120.045 s & Unfinished.\\
Largest target, structural rational form & \code{ByteCount} & 3,517,376 & 6,812,872 & 50.350 s & Expanded.\\
Common-factor extraction & \code{ByteCount} & 1,362,792 & 684,744 & 3.393 s & Fraction-free targets.\\
Factored-route final cleanup & \code{ByteCount} & 855,104 & 759,840 & 10.741 s & Full factor-first route about 37.76 s.\\
Second master sweep & \code{ByteCount} & 759,840 & 759,840 & 9.426 s & No change.\\
Direct target cleanup & \code{ByteCount} & 1,542,536 & 727,376 & 450.188 s & Much slower than factor-first.\\
33-rule certified normalization & \code{ByteCount} & 1,035,984 & 283,144 & 9.437 s & Exact fraction independence.\\
Divide only & \code{ByteCount} & 1,035,984 & 1,089,000 & 0.800 s & Fractions remain.\\
Divide then \code{Simplify} & \code{ByteCount} & 1,035,984 & 283,144 & 5.877 s & Best compact strip variant.\\
Divide then \code{FactorTerms} & \code{ByteCount} & 1,035,984 & 1,090,008 & 5.384 s & Fractions remain.\\
Final \code{FullSimplify} & \code{ByteCount} & 283,144 & 256,456 & 32.352 s & Modest extra reduction.\\
Final \code{Factor} & \code{ByteCount} & 283,144 & 925,968 & 0.146 s & Expanded.\\
Final color collection & \code{ByteCount} & 283,144 & 562,496 & 0.185 s & Expanded.\\
Signature cleanup & \code{ByteCount} & 1,035,984 & 3,879,816 & 1.868 s & Expanded 3.745-fold.\\
Strip only & \code{ByteCount} & 1,035,984 & 1,935,248 & 1.741 s & Expanded.\\
Term cleanup & \code{ByteCount} & 1,035,984 & 1,935,248 & 1.691 s & Expanded.\\
Denominator groups & \code{ByteCount} & 1,035,984 & 1,937,768 & 1.905 s & Expanded.\\
Color collection variant & \code{ByteCount} & 1,035,984 & 1,935,248 & 1.821 s & Expanded.\\
Automatic common factor & Syntactic factor & -- & incomplete & 0.0067 s & Missed $1/(x_a x_b)$.\\
\bottomrule
\end{longtable}
\scriptsize
\begin{longtable}{@{}llrrrl@{}}
\caption{NLO TT and NNLO simplification measurements.}\label{tab:all-ttnnlo}\\
\toprule
Method & Object and metric & Input & Output & Time & Interpretation \\
\midrule
\endfirsthead
\toprule
Method & Object and metric & Input & Output & Time & Interpretation \\
\midrule
\endhead
TT incorrect tensor route & \code{ByteCount} & -- & 184,796,536 & -- & Wrong dimension replacement; 39,036 GLIs.\\
TT corrected sparse cut route & \code{ByteCount} & -- & 27,754,864 & -- & Exact tensor tests; nine GLIs.\\
TT target first, 5 by 5 & \code{ByteCount} & 139,298,664 & 10,562,384 & 52.02 s & Four kernels.\\
TT final master cleanup & serialized bytes & -- & 750,504 & 30.36 s & Six masters.\\
TT angular master & \code{ByteCount} & 1,224,256 & 119,520 & 1.874 s & Two angular coefficients.\\
64 largest NNLO targets & \code{ByteCount} & 83,386,848 & 28,279,424 & 75.643 s & Four kernels; 2.949-fold.\\
Largest target, whole \code{Simplify} & \code{ByteCount} & 2,434,656 & 811,576 & 15.30 s & Exact zero difference.\\
Largest target, termwise \code{Simplify} & \code{ByteCount} & 2,434,656 & 701,328 & 6.92 s & Exact zero difference.\\
Largest target, \code{Cancel[Together]} & \code{ByteCount} & 2,434,656 & 3,748,912 & 0.37 s & Expanded.\\
Largest target, add \code{FactorTerms} & \code{ByteCount} & 2,434,656 & 5,316,312 & 0.83 s & Expanded further.\\
Whole-master \code{Simplify} & serialized bytes & 4.7--129.8 MB & -- & 1,800 s each & Eight jobs unfinished; about 34 GiB total.\\
Master 8 termwise & serialized bytes & 4,655,210 & 3,082,976 & 175.070 s & 377 terms, eight kernels.\\
Master 1 signatures & serialized bytes & 1,607,146,927 & 920,945,177 & about 19.5 min & Rational normalization within signatures.\\
Master 2 signatures & serialized bytes & 903,734,615 & 464,281,765 & about 15 min & Rational normalization within signatures.\\
Master 4 signatures & serialized bytes & 83,696,010 & 47,987,161 & 68.55 s & 2,591 different signatures.\\
Old all-master cleanup & serialized bytes & 3,960,695,102 & 3,597,827,125 & -- & Only 1.101-fold.\\
Symbolica then Mathematica & serialized MB & -- & 0.81 MB & 12.16 s & Direct Mathematica: 0.70 MB in 5.91 s.\\
125-target root sample & \code{ByteCount} & 52,630,888 & 559,984 & 30.45 s & Four kernels; certified root grammar.\\
Entrywise root cleanup & serialized bytes & 859,107,253 & 131,787,741 & 1,304.639 s & Three post-IBP master columns, eight kernels.\\
Eager denominator merge & \code{ByteCount} & 422,282,520 & 1,727,567,608 & 6,375.22 s & 6,290 fractions; expanded strongly.\\
Size-monotone merge & \code{ByteCount} & 422,282,520 & 351,559,840 & 1,260.216 s & 6,557 fractions; serialized 99,892,985 bytes.\\
Dimensionless scale separation & \code{ByteCount} & 351,559,840 & 296,802,344 & 430.914 s & Serialized 99,892,985 to 92,385,524.\\
Dimensionless denominator merge & \code{ByteCount} & 296,802,344 & 296,275,768 & 904.575 s & 5,973 fractions; serialized 92,241,826.\\
Denominator-factor inventory & serialized bytes & 92,241,826 & 91,385,453 & 27.94 s & Only 0.93 percent smaller; 30 factors.\\
\bottomrule
\end{longtable}
\end{landscape}

\section{Provenance and retained records}

The documentation directory contains:
\begin{itemize}
  \item the two earlier narrative records;
  \item the compact NLO UU, TT, and NNLO manifests quoted above;
  \item scripts for every principal benchmark and branch-safe transformation;
  \item a machine-readable CSV table naming every size metric;
  \item an artifact index pointing to the original large expressions and their
  stored hashes.
\end{itemize}
The report does not duplicate multi-gigabyte coefficient files.

\end{document}
